\chapter{AMP, TAP, and the Bulk Fixed Point}\label{app:amp}

This appendix develops the per-node Gaussian closures of statistical
physics --- AMP, the TAP equations --- on the Gaussian chain, and proves
the exact breakdown announced in Remark~\ref{rmk:bp-no-clt}. The results
are complete and audited; they are collected here rather than in
Chapter~\ref{ch:bp} because they sit off the thesis's main line. The
punchline: on the chain, AMP still returns the \emph{exact score}, but its
variance closure develops an exact breakdown time $t_c(\alpha)$ with
critical coupling $\alpha_c = \sqrt{2} - 1$ --- the supervision remark
that ``a central limit theorem on two data points is usually not right'',
turned into a theorem with a phase diagram.

\section{The bulk cavity fixed point of BP}\label{sec:amp-bulk-bp}

Chapter~\ref{ch:gaussian} solved the bulk covariance by matrix algebra;
the same formulas re-emerge from the message recursion, which is worth
exhibiting because it localises \emph{where} BP and AMP differ. In the
bulk, parametrise the forward message (evidence absorbed) by its precision
$\lambda_k$; the closure lemmas of Section~\ref{sec:bp-gaussian} give the
scalar recursion
\begin{equation}
  \lambda_{k+1} \;=\; J_d - \frac{\beta^2}{\lambda_k},
  \label{eq:bp-cavity}
\end{equation}
with $J_d, \beta$ the bulk parameters \eqref{eq:g-bulk-params}: absorbing
one neighbour's summary costs the Gaussian-integration correction
$-\beta^2/\lambda$. The fixed points are
$\lambda^\pm = \big( J_d \pm \sqrt{J_d^2 - 4\beta^2} \big)/2$, and since
$J_d - 2|\beta| = e^{-2t}/\Deltat + (1 - |\alpha|)^2/\sigma_\eta^2 > 0$
(the AM--GM margin of Chapter~\ref{ch:gaussian}), the attractive upper
fixed point
\begin{equation}
  \boxed{\;\lambda^\ast = \frac{J_d + \sqrt{J_d^2 - 4\beta^2}}{2}\;}
  \label{eq:bp-lambdastar}
\end{equation}
exists for \emph{every} $\alpha \in (-1,1)$ and $t > 0$, with
$|f'(\lambda^\ast)| = \beta^2/\lambda^{\ast 2} < 1$: the sweep converges
geometrically from any initialisation. Recombining the two directions and
the local term, and using the fixed-point relation
$\beta^2/\lambda^\ast = J_d - \lambda^\ast$,
\begin{equation}
  \Lambda_{\mathrm{bulk}}
  = J_d - \frac{2\beta^2}{\lambda^\ast}
  = 2\lambda^\ast - J_d
  = \sqrt{J_d^2 - 4\beta^2}
  \quad\Longrightarrow\quad
  V_{\mathrm{exact}} = \frac{1}{\sqrt{J_d^2 - 4\beta^2}} ,
  \label{eq:bp-bulkvar}
\end{equation}
in exact agreement with Theorem~\ref{thm:g-bulk-cov}. \emph{Belief
propagation on the Gaussian chain has no phase transition}: the cavity
fixed point exists always, and the discriminant $J_d^2 - 4\beta^2$ is the
quantity to watch as we now modify the closure.

\section{AMP/TAP: same score, different variance}\label{sec:bp-amp}

AMP --- the TAP equations of spin-glass theory
\citep{thouless1977solution, donoho2009message, zdeborova2016statistical}
--- is the message passing one uses when the factor graph is \emph{dense}:
each message is then a sum of many weak contributions, approximately
Gaussian by a central-limit argument, and one tracks per-\emph{node} means
and variances instead of per-\emph{edge} messages, with the Onsager term
correcting the self-influence. Rigorous derivations of AMP from BP exist
in exactly that dense limit \citep{genovese2026amp}. Our chain is the
opposite regime in every respect --- diagonal measurement, coupling in the
prior, degree two --- so it is the cleanest possible place to see what the
per-node closure costs. Viewing the posterior as a Gaussian MRF with
precision $J$ and field $h$, the three schemes solve, at each node,
\begin{equation}
  \text{mean field: } V_i = \frac{1}{J_{ii}};
  \qquad
  \text{AMP/TAP: } V_i = \frac{1}{J_{ii} - \sum_{k \neq i} J_{ik}^2 V_k};
  \qquad
  \text{BP: } V_i = (J^{-1})_{ii} .
  \label{eq:bp-three-closures}
\end{equation}
The AMP closure is the BP cavity with the ``exclude the receiving
neighbour'' correction dropped ($V_{k \to i} \approx V_k$): an $O(1/N)$
change on a dense graph, an $O(1)$ change on a chain with two neighbours.

\begin{theorem}[The score is closure-independent]\label{thm:bp-same-score}
For a Gaussian posterior, the fixed-point condition for the \emph{means}
is the same linear system $J m = h$ under mean field, BP, and AMP. All
three schemes hence return the exact posterior mean and --- via the
score--posterior identity --- the same, exact, joint score, for every
$(K, \alpha, t)$, regardless of whether their variance closures are
accurate or even solvable.
\end{theorem}

The proof is one line: in each scheme the mean update at node $i$ reads
$m_i \leftarrow (h_i - \sum_{k \neq i} J_{ik} m_k)/J_{ii}$ at the fixed
point --- the closures modify the uncertainty bookkeeping, not the linear
relation among means --- and the unique solution of $Jm = h$ is the exact
posterior mean. Numerically: AMP means agree with exact means to
$4.3 \times 10^{-12}$, and AMP scores with BP scores to
$1.7 \times 10^{-12}$, across all tested chain configurations.

For the \emph{variances}, the story splits. In the bulk the AMP closure
\eqref{eq:bp-three-closures} reads $V = 1/(J_d - 2\beta^2 V)$ --- compare
BP's cavity $V_{\mathrm{cav}} = 1/(J_d - \beta^2 V_{\mathrm{cav}})$:
\emph{the entire difference between BP and AMP on the chain is a factor of
two}, both neighbours included where BP excludes one.

\begin{theorem}[AMP bulk variance, breakdown time, critical coupling]
\label{thm:bp-amp-breakdown}
The AMP bulk closure is the quadratic $2\beta^2 V^2 - J_d V + 1 = 0$, with
physical root
\begin{equation}
  V_{\mathrm{AMP}}
  = \frac{J_d - \sqrt{J_d^2 - 8\beta^2}}{4\beta^2},
  \label{eq:bp-amp-var}
\end{equation}
which exists iff $J_d \geq 2\sqrt{2}\,|\beta|$ --- strictly stronger than
BP's always-true $J_d > 2|\beta|$. Since $J_d(t)$ decreases monotonically
in $t$, the existence condition fails, if ever, past a single time;
solving it as an equality (with $u := e^{-2t}$,
$e^{-2t}/\Deltat = u/(1-u)$) gives
\begin{equation}
  \boxed{\;
  t_c(\alpha) = -\tfrac{1}{2} \log \frac{g}{1+g},
  \qquad
  g(\alpha) = \frac{2\sqrt{2}\,|\alpha| - 1 - \alpha^2}{1 - \alpha^2},
  \;}
  \label{eq:bp-tc}
\end{equation}
finite exactly when $g > 0$, i.e.
\begin{equation}
  \boxed{\;|\alpha| \;>\; \alpha_c \;=\; \sqrt{2} - 1 \;\approx\; 0.4142 .\;}
  \label{eq:bp-alphac}
\end{equation}
For $|\alpha| \leq \alpha_c$ the AMP variance has a fixed point at every
diffusion time; for $|\alpha| > \alpha_c$ it ceases to exist for all
$t > t_c(\alpha)$. In the weak-coupling regime the two closures agree to
second order, with
\begin{equation}
  V_{\mathrm{AMP}} - V_{\mathrm{exact}}
  = \frac{2\beta^4}{J_d^5} + O\!\Big(\frac{\beta^6}{J_d^7}\Big)
  \;>\; 0 :
  \label{eq:bp-weak}
\end{equation}
AMP always \emph{over}estimates the bulk variance, at fourth order in the
coupling.
\end{theorem}

\begin{toolbox}{Roots, branch, and the phase boundary}
\emph{Branch selection.} The quadratic has roots
$V_\pm = (J_d \pm \sqrt{J_d^2 - 8\beta^2})/(4\beta^2)$; as $\beta \to 0$,
expanding the square root gives $V_- \to 1/J_d$ (the decoupled answer)
while $V_+$ diverges --- $V_-$ is the physical branch, and the attractive
fixed point of the damped iteration used in practice.
\emph{Critical coupling.} $g > 0 \iff \alpha^2 - 2\sqrt{2}|\alpha| + 1 < 0
\iff |\alpha| \in (\sqrt{2}-1, \sqrt{2}+1)$; intersecting with
$|\alpha| < 1$ leaves $|\alpha| > \sqrt{2} - 1$. Sanity check at
$t = \infty$: there the posterior is the bare prior, and the existence
condition applied to $\Qz$ itself reads $(1 + \alpha^2) \geq
2\sqrt{2}|\alpha|$ --- AMP applied to a bare AR(1) chain works iff
$|\alpha| \leq \sqrt{2} - 1$, which is \eqref{eq:bp-alphac} read
backwards. \emph{Weak coupling.} With $\epsilon := \beta^2/J_d^2$:
$(1 - 4\epsilon)^{-1/2} = 1 + 2\epsilon + 6\epsilon^2 + O(\epsilon^3)$
for the exact variance;
$1 - \sqrt{1 - 8\epsilon} = 4\epsilon + 8\epsilon^2 + O(\epsilon^3)$,
divided by $4\epsilon J_d$, gives
$V_{\mathrm{AMP}} = (1 + 2\epsilon + 8\epsilon^2 + O(\epsilon^3))/J_d$;
subtract. \hfill$\square$
\end{toolbox}

\begin{figure}[t]\centering
\includegraphics[width=\linewidth]{figures/fig_bulk_variance.pdf}
\caption{\textbf{The three closures in closed form.} Left
($\alpha = 0.8 > \alpha_c$): exact $=$ BP, AMP \eqref{eq:bp-amp-var}, and
mean field; the AMP branch terminates with a square-root singularity at
the exact breakdown time $t_c(0.8) = 0.228$ from \eqref{eq:bp-tc}
(shaded: no fixed point). Right ($\alpha = 0.3 < \alpha_c$): AMP exists at
every $t$ and slightly overestimates the variance, as predicted by
\eqref{eq:bp-weak}.}
\label{fig:bp-bulkvar}
\end{figure}

\begin{figure}[t]\centering
\includegraphics[width=\linewidth]{figures/fig_bp_vs_amp.pdf}
\caption{\textbf{BP vs AMP on a finite chain} ($K = 9$, $\alpha = 0.8$).
Left: the AMP \emph{mean} error ($\sim 10^{-12}$, flat in $t$: the score
is closure-independent, Theorem~\ref{thm:bp-same-score}) against the AMP
\emph{variance} error, which grows with $t$ until the closure stops
converging (shaded). Right: existence of the AMP variance fixed point in
the $(\alpha, t)$ plane; the black curve is the exact boundary
$t_c(\alpha)$ of \eqref{eq:bp-tc}, the dashed line
$\alpha_c = \sqrt{2}-1$; theory and iteration agree at all $180$ scanned
points.}
\label{fig:bp-vs-amp}
\end{figure}

Numerically, on a $K = 12$ chain at $\alpha = 0.8$: the AMP score error
stays at $10^{-13}$--$10^{-12}$ for every $t$, while the variance error is
$1.3 \times 10^{-4}$ at $t = 0.05$, $3.5 \times 10^{-2}$ at $t = 0.2$, and
past $t \approx 0.23$ the variance update has no positive fixed point at
all. The audit further verifies: the closed form \eqref{eq:bp-amp-var}
against the converged iteration to $2.5 \times 10^{-12}$; the existence
condition at all $180$ points of an $(\alpha, t)$ scan; $t_c$ against the
bisected flip point of the iteration to $6.6 \times 10^{-6}$ relative;
convergence below $\alpha_c$ even at $t = 50$; and the weak-coupling ratio
$(V_{\mathrm{AMP}} - V_{\mathrm{exact}})/(2\beta^4/J_d^5) \to 1$ within
$0.4\%$.

The supervision remark that framed this appendix --- the Gaussian message
assumption is a CLT statement, and ``a central limit theorem on two data
points is usually not right'' --- is thus made exact, with both halves
landing precisely where predicted. The \emph{conjecture} (Gaussian message
passing on the solved case recovers the known score) holds exactly: for BP
by algebraic closure, for AMP because the mean fixed point is
closure-independent. The \emph{objection} targets the variance closure,
and there it is not merely right but quantitative: a factor two in a
discriminant, a phase boundary \eqref{eq:bp-tc}, and a critical coupling
$\sqrt{2} - 1$. AMP is not wrong --- on a dense benchmark (random
$N = 600$ SPD precision) the same closure gets variances to
$1.7 \times 10^{-3}$ where mean field errs at $4 \times 10^{-2}$ --- it is
built for a different graph.
